% Separate research fragment, 7 October 2026. No frozen source, accepted root,
% open manuscript or compiler state is changed by this fragment.
% Analytical proofs are distinguished from Lean acceptance and conjectures.

\subsection{Expansion preserves the densest core selections}

For a finite $s$-uniform hypergraph $F$ with $v(F)>s$, write
\[
 m_s(F)=\max_{Q\subseteq F,\ v(Q)>s}
           \frac{e(Q)-1}{v(Q)-s}.
\]
Vertices not incident with a selected edge are permitted in $Q$.
Here balancedness means that the whole eligible graph attains this maximum;
strictness additionally excludes every proper eligible maximizer. The usual
definition in the cited expansion paper is restricted to at least three
edges; our explicit two-edge and one-edge discussion uses this literal
density-comparison extension, not a change to that paper's convention.
For a graph $G$, its private $r$-expansion replaces every edge $uv$ by
$\{u,v\}\cup P_{uv}$, where the sets $P_{uv}$ have size $r-2$, are pairwise
disjoint, and avoid the original vertices. Original isolated vertices are retained.

\begin{proposition}[The known expansion transform, with equality cases]
Let $G$ be a finite simple graph with at least two edges, and let $r\ge2$.
Then
\[
 m_r(G^{[r]})
   =\frac{m_2(G)}{1+(r-2)m_2(G)},\qquad
 \frac1{m_r(G^{[r]})}=r-2+\frac1{m_2(G)}.
\]
The graph is $2$-balanced if and only if its expansion is $r$-balanced,
and the corresponding strict-balancedness conditions are equivalent.
More precisely, an eligible selected expanded subgraph attains $m_r$
if and only if it has at least two edges, contains no vertex outside their
union, and its selected core edges with their incident vertices attain $m_2$.
\end{proposition}

The transform itself is already Nie--Spiro's Lemma~4.2, which is stated
more generally for expansions of uniform hypergraphs
\citep{niespiro_expansions2024}. We use it as an attributed structural
bridge, not as a new result.

\begin{proof}
Two distinct graph edges have three or four incident vertices, so $m_2(G)>0$.
A maximizer therefore has at least two edges and no isolated selected vertex:
zero- and one-edge selections have nonpositive density, while deleting an
isolated selected vertex strictly increases positive density.
For a set $S$ of $f\ge2$ core edges, let $v$ be the number of its incident
vertices. Its expanded edge union has exactly $v+(r-2)f$ vertices. Hence
\[
 \frac{f-1}{v+(r-2)f-r}
   =T_r\left(\frac{f-1}{v-2}\right),\qquad
 T_r(x)=\frac{x}{1+(r-2)x}.
\]
All denominators are positive, and $T_r$ is strictly increasing for $x>0$.
Every expanded edge selection projects to one core edge selection, and extra
selected vertices strictly lower its positive density. Taking maxima gives
the identity and the stated correspondence of equality cases.

The whole-graph ratio transforms by the same formula, proving balancedness
equivalence. A balanced graph with at least two edges cannot have isolated
vertices. Thus a proper expanded maximizer corresponds to a proper core
maximizer, and conversely a proper core maximizer expands to a proper one.
This proves strictness equivalence, including the two-edge cases.
\end{proof}

The domain restrictions matter. A single edge with at least one isolated
vertex has $m_2=m_r=0$; it is strictly balanced exactly when there is one
isolated vertex. An edgeless graph on $n>r$ vertices has the literal values
$m_2=-1/(n-2)$ and $m_r=-1/(n-r)$, so the algebraic identity also holds there;
these negative values are not used as positive random extremal parameters.
If the eligible-subgraph domain is empty, its maximum is undefined here.
In particular $K_2$ and its expansion have no eligible subgraph, and an
edgeless graph with $3\le n\le r$ has a defined $m_2$ but no defined $m_r$.
No arbitrary zero convention is inserted into the theorem.

For an edge-containing forest with at least two edges, $m_2=1$ if it has
two adjacent edges, and $m_2=1/2$ if it is a matching. The corresponding
expanded values are $1/(r-1)$ and $1/r$. A connected tree with at least
three edges is balanced but not strictly balanced, because a proper
two-edge path is already a maximizer.

\subsection{A family certified by spanning triangular exposure orders}

Let $N\ge1$ and $t\ge2$. Consider $tN$ distinct face labels, partitioned
into $t$ ordered classes of $N$ labels each, and a common set of $N+2$
point labels. Every face contains exactly three distinct points. In every
class, each face after the first has a distinguished point absent from all
earlier faces. Repeated point-triples in different face classes are allowed;
the ordinary simple triangular-complex case is a special case. Face labels
remain distinct vertices in the bipartite incidence graph $J$.

\begin{proposition}[Spanning exposure certificate]
Under these hypotheses, $J$ is $2$-vertex-connected and strictly $2$-balanced.
Its exact parameters are
\[
 e(J)=3tN,\quad v(J)=(t+1)N+2,\quad
 m_2(J)=\frac{3tN-1}{(t+1)N},\quad
 \beta(J)=(2t-1)N-1.
\]
Consequently its private $r$-expansion is strictly $r$-balanced with
\[
 m_r(J^{[r]})=
 \frac{3tN-1}{(t+1)N+(r-2)(3tN-1)}.
\]
For $r\ge3$, this gives a bipartite $2$-connected core admissible as a
single obligatory atom under the higher-uniformity classification.
The certificate says nothing by itself about Sidorenko's inequality.
\end{proposition}

\begin{proof}
The first face of a class uses three points, and its other $N-1$ faces
each introduce a fresh point. Since the entire ambient point set has size
$N+2$, each class spans all points and every later face introduces exactly
one point while meeting the earlier union in exactly two points.
After deleting any point, this ordered incidence construction stays connected:
the first face retains at least two points, and every later face retains an
old point through which it attaches. After deleting any face label, another
complete class connects all point labels; every remaining face attaches to
those points. Thus deleting any vertex leaves a connected graph.

Consider any selected incidence subgraph with $a$ point nodes, $b$ face
nodes and $e$ edges. If $a\ge2$, each class contributes at most $a-2$
complete faces: the earliest of any $h\ge1$ retained faces supplies three
points and every later retained face supplies its fresh point, so their
union has at least $h+2$ points. Hence
\[
 e\le\min\{3b,\ 2b+t(a-2)\}.
\]
Also $e\le ab$, $a\le N+2$ and $b\le tN$. Selections with $a=0$ have
negative density, those with $a=1$ have density at most one, and those
with $a=2$ have density at most $2-1/b$. All are strictly below the
whole-graph value $D_N=(3tN-1)/((t+1)N)$ whenever their denominator is positive.
For $t=2$, use $b\le2N$; for $t\ge3$, use $D_N\ge2$.

It remains to consider $a\ge3$. Put $x=a-2$, so $1\le x\le N$, and
$D_x=(3tx-1)/((t+1)x)$. If $b\le tx$, then
\[
 \frac{e-1}{x+b}\le\frac{3b-1}{x+b}\le D_x\le D_N,
\]
because the middle fraction is strictly increasing in $b$ and $D_x$
is strictly increasing in $x$. If $b\ge tx$ and $t\ge3$, then
$(2b+tx-1)/(x+b)$ is non-increasing in $b$: its derivative has sign
$1-(t-2)x\le0$. The same upper bound $D_x\le D_N$ follows.
If $t=2$, the remaining expression is $2-1/(x+b)$; the bounds
$x\le N$ and $b\le2N$ give $2-1/(x+b)\le2-1/(3N)=D_N$.

Equality in either route forces $x=N$ and $b=tN$; it also forces
$e=3tN$. This is the full incidence graph. The apparent plateau
$t=3,x=1$ creates no extra equality case: equality with $D_N$ first
forces $N=1$, after which $b\ge t$ and $b\le tN$ force $b=t$.
Every proper eligible subgraph is therefore strict. The exact counts
and connected cycle rank follow, and the expansion claim uses the
preceding transform.
\end{proof}

For $t=2,N=11$, the parameters are $35$ vertices, $66$ edges and rank $32$,
recovering the incidence example already under study. The case $N=1$
is $K_{3,t}$ when repeated point-triples are allowed. For a simple
triangular complex, $N=2,t=2$ is realized by the four triples on four
points and gives the cube graph, with $m_2=11/6$ and cycle rank five.
Neither example is offered as a new Sidorenko theorem.

The spanning-size hypothesis is sharp for this route. On five points,
take two classes both ordered as $abc,ade$. They have the same freshness
property with $N=2$, and both span all five points, but the common point
$a$ is a cut vertex. The whole incidence graph has density $11/7$,
whereas the two copies of $abc$ on their three points have density $5/3$.
Thus a freshness certificate alone, without the exact $N+2$ saturation,
does not imply balancedness or the required connectivity.

\subsection{The unconditional random-host benchmark}

Let $F$ be a fixed strictly balanced private expansion from the preceding
certificate, and put $\alpha=1/m_r(F)$ and $\gamma=r-\alpha$. Then
\[
 \alpha=r-2+\frac{(t+1)N}{3tN-1},\qquad
 \gamma=2-\frac{(t+1)N}{3tN-1}.
\]
Thus the deletion-scale exponent $\gamma$ is independent of uniformity.
This is a random Tur\'an deletion scale, not the first-appearance threshold.
It does not require positive Sidorenko excess.

If $\rho n^\alpha\to0$ and $\rho n^r\to\infty$, then asymptotically
almost surely
\[
 \operatorname{ex}(G^{(r)}_{n,\rho},F)
       =(1-o(1))\,\rho\binom nr.
\]
Indeed, the expected number of copies divided by the expected edge count
is $O((\rho n^\alpha)^{e(F)-1})=o(1)$. Markov's inequality and edge-count
concentration show that removing one edge from every remaining copy loses
only an $o(1)$ fraction of the host edges. If $\rho\ge n^{-\alpha}$,
independent thinning to probability $n^{-\alpha}/\log n$ gives instead
\[
 \operatorname{ex}(G^{(r)}_{n,\rho},F)\ge n^{\gamma-o(1)}
 \quad\text{asymptotically almost surely}.
\]
These elementary benchmark bounds must not be confused with a polynomial
gain from non-Sidorenko excess or with a matching upper bound.

\subsection{Quantitative extraction and an equally visible conjectural direction}

\paragraph{Source-backed extraction and validation target.}
Let $J$ be the fixed $35$-vertex, $66$-edge incidence core, and let
$A_r=J^{[r]}$. There exist constants $c,C>0$ and $r_0$ such that, for every
integer $r\ge r_0$, there is a finite uniform probability space $\Omega_r$
with $|\Omega_r|\le r^C$ and a symmetric, one-coordinate-regular kernel
$U_r:\Omega_r^r\to[0,1]$ of mean $p_r\ge r^{-Cr}$ satisfying
\[
t(A_r,U_r)\le(1-c r^{-66})p_r^{66}.
\]
Here $t(F,U)$ is the expectation of the product of $U$ over the edges of
$F$, with independent uniform samples at all vertices, including all
noninjective assignments. One-coordinate regularity means that integrating
all but any one coordinate gives the same constant $p_r$.
This is a quantitative target, not a statement established by the currently
audited rank-layer transfer, and its originality has not been certified.
In particular, it is not advertised as a newly identified open conjecture:
the source contains the power savings needed for a quantitative extraction.
It would imply $s(A_r)\ge c'/(r^{67}\log r)$ after the usual kernel-to-finite-host
approximation, and hence a quantified random-host gain.

The exact unverified extraction is a polynomial-rate proper-subgraph-flat
family: symmetric regular nonnegative pair kernels $K_q$ of mean one with
$|\Omega_q|,\|K_q\|_\infty\le C_0q^B$, every proper edge-subgraph moment
within $C_0q^{-\eta}$ of one, and the full $J$-moment at most $1-c_0$,
for constants $B,\eta,c_0>0$ and every sufficiently large odd prime $q$.
The primary source does contain power savings: its transverse estimates are
$O_D(q^{-1/2})$, and the tail estimates support a fixed-dimension
$O_D(q^{-1/6})$ rank-layer extraction, although its final statement uses
$o(1)$. Bounds polynomial in $q$ for the alphabet and kernel supremum are
likewise extractable for fixed $D$~\citep{openai161}. Proper-edge flatness
and exact row/column regularity have separate source-derived checks.
Thus a lack of source rates is not the obstruction. What remains unproved
as a complete synthesis in this bounded audit is the full uniform-space
family above, including rational type-prior replication and validation of
the external rank-layer input. Those tasks are not discharged by reading
an asymptotic headline alone.

Given that input, choose an odd prime $q$ polynomially large in $r$ so the moment error
is at most $(c_0/2)\lambda_r^{66}$, where $\lambda_r\asymp1/r$ is the matching
profile coefficient. The exact edge-subset expansion gives a relative
deficit at least $(c_0/2)\lambda_r^{66}$. The finite matching extension is
nonnegative and regular, with maximum at most
$\max\{1,\|K_q\|_\infty\}^{\lfloor r/2\rfloor}$; normalization therefore
gives $p_r\ge r^{-Cr}$. This is a complete conditional route, not a substitute
for proving the stated polynomial-rate input.

For completeness, the matching extension averages
$\prod_{ij\in M}K_q(X_i,X_j)$ over all maximum matchings $M$ of the $r$
coordinates. Each term has mean one and constant one-coordinate marginals,
because $K_q$ is regular of mean one. A fixed coordinate pair is matched
with probability $1/(r-1)$ for even $r$ and $1/r$ for odd $r$; integrating
the other coordinates therefore gives $1+\lambda_r(K_q-1)$. Distinct
private ports in $A_r$ then give $t(A_r,W_r)=t(J,1+\lambda_r(K_q-1))$.
Expanding the product over the $66$ core edges uses nonnegative coefficients
$\lambda_r^{|S|}(1-\lambda_r)^{66-|S|}$ summing to one. Proper-subgraph
errors cost at most $\varepsilon(1-\lambda_r^{66})$, whereas the full term
gains $c_0\lambda_r^{66}$. The chosen error makes the net gain at least
$(c_0/2)\lambda_r^{66}$. If the extension is normalized by an upper bound
$B_r>1$, then $p_r=B_r^{-1}$ and the relative deficit is unchanged.
Finally $-\log(1-\delta)\ge\delta$ and $-\log p_r\le Cr\log r$ give the
stated excess lower bound. These are finite algebraic steps; the analytic
rate input and kernel-to-finite-host conversion remain separate.

A simpler one-pair average avoids the matching extension: use
$W_r=\binom r2^{-1}\sum_{i<j}K_q(X_i,X_j)$. Its pair marginal is
$1+\binom r2^{-1}(K_q-1)$ and its maximum is polynomially bounded.
Under the same rate input, it gives the weaker relative deficit
$c r^{-132}$ but the stronger normalization $p_r\ge r^{-C}$.
This trade-off identifies a concrete question about witness complexity,
rather than another unqualified existence assertion.

\paragraph{Working conjectural direction: remove the exponential normalization loss.}
For the same fixed $J$ and $A_r$, do there exist $c,C>0$ and $r_0$ such
that every integer $r\ge r_0$ admits a uniform space $\Omega_r$ with
$|\Omega_r|\le r^C$ and a symmetric, one-coordinate-regular
$U_r:\Omega_r^r\to[0,1]$ satisfying
\[
 p_r\ge r^{-C},\qquad
 t(A_r,U_r)\le(1-c r^{-66})p_r^{66}?
\]
This stronger target combines the polynomial mean of the one-pair lift
with the larger deficit of the matching lift. Neither current route proves
it. If true, it would improve the excess bound to
$s(A_r)\ge c'/(r^{66}\log r)$. No novelty, established-openness or sharpness
claim is attached to this working question.

A proof would need sharper normalization control for a regular lift or a
different witness, rather than merely reusing the matching supremum bound.
Failure of that particular construction is not a counterexample to the
question. One sufficient disproof would be a uniform trade-off bound: for
every fixed $C$, along an unbounded sequence of $r$ the supremum of
\[
 \left(1-\frac{t(A_r,U)}{p(U)^{66}}\right)_+
\]
over all such uniform regular kernels with $|\Omega|\le r^C$ and
$p(U)\ge r^{-C}$ is $o(r^{-66})$. This makes the proof and falsification
obligations exact without presuming that the desired improvement exists.

% Formal status: the transform and this follow-on theorem are analytical;
% no new Lean source or compiler attempt is included in this fragment.
% Frozen candidate counters and the accepted root remain unchanged.
